% \yc{suggestion: add a preamble here. summarize the key findings and the results here?}
\subsubsection{Main Results}
\label{sec:main results}

Here we present the results for our base models and several other well-known base models across standard academic benchmarks. While benchmarking open-source models, we observed a disparity between the results generated by our pipeline and those reported in public sources. Upon conducting a more in-depth investigation of this difference, mostly because different models use different prompts, post-processing strategies, and sampling techniques.
These differences may potentially induce significant variations in the outcomes. Our prompt and post-processing strategy remains consistent with the default settings of the original benchmarks\citep{PIQA,siqa,hellaswag,winogrande,arc,obqa,csqa,squad,quac,boolq,cobbe2021training,hendrycks2021measuring,chen2021evaluating,austin2021program,hendrycks2020measuring,li2023cmmlu,zhang2023evaluating,srivastava2023imitation,suzgun2022challenging}.
We use greedy decoding without any post-processing for the generated content. For scores that were not reported publicly (or scores reported with different settings), we try to get results with our pipeline. For scores that can be found publicly, we directly report the existing numbers. We use the following benchmarks, largely following the practice of LLaMA 2~\citep{llama2}:

\begin{description}[leftmargin=0.2in]
    \item[Commonsense Reasoning:] We included PIQA\cite{PIQA}, SIQA\cite{siqa}, HellaSwag\cite{hellaswag}, WinoGrande \cite{winogrande}, ARC\cite{arc}, OpenBookQA(OBQA)\cite{obqa}, and CommonsenseQA(CSQA)\cite{csqa} to assess common sense reasoning. CSQA was exclusively tested using a 7-shot setup, while all other tests were conducted with a 0-shot configuration.
    \item[Reading Comprehension:] For reading comprehension, we report the 0-shot average on SQuAD\cite{squad}, QuAC\cite{quac}, and BoolQ\cite{boolq}.
    \item[Math:] We report the average of the GSM8K\cite{cobbe2021training} (8 shot), and MATH\cite{hendrycks2021measuring} (4 shot) benchmarks with pass@1 accuracy without any specific prompting strategy (e.g. Chain-of-Thought prompting) and other ensemble technique (e.g., majority voting).
    \item[Code:] We report the average pass@1 scores of our models on HumanEval\cite{chen2021evaluating}  (Chen et al., 2021) and MBPP\cite{austin2021program} (Austin et al., 2021). 
    \item[Popular Aggregated Benchmark:] We report the overall results for MMLU\cite{hendrycks2020measuring}(5-shot), CMMLU\cite{li2023cmmlu} (5-shot), Gaokao-Bench\cite{zhang2023evaluating} (5-shot), and BigBench\cite{srivastava2023imitation} Hard (BBH\cite{suzgun2022challenging}) (3-shot).
\end{description}


% \begin{table}
%   \centering
%   \begin{tabular}{lcccc}
% \toprule
%               \bf Model & \bf Size & \bf Avg & \bf GSM8k & \bf MATH \\
% \midrule
% \bf GPT-3.5 & - & 35.6 & 57.1 & 14.0 \\
% \bf GPT-4 & - & \textbf{66.1} & \textbf{92} & \textbf{40.2} \\
% \gmidrule
% \bf Falcon & 180B & - & 54.4 & - \\
% \gmidrule
% \multirow{2}{*}{\bf Qwen}
%  & 7B & 31.7 & 51.7 & 11.6 \\
% & 14B & 43.1 & 61.3 & 24.8 \\
%  \gmidrule
% \multirow{2}{*}{\bf Baichuan 2}
% & 7B & 15.1 & 24.5 & 5.6 \\
% & 13B & 16.1 & 22.1 & 10.1 \\
% % \midrule
% % Open-L\textsmaller{LaMA} v2 & 7B & 9.7 & 2.1 \\
% \gmidrule
% \multirow{3}{*}{\bf L\textsmaller{LaMA} 2} 
% & 7B & 10.0 & 16.7 & 3.3\\ 
% & 34B & 24.2 & 42.2 & 6.2\\ 
% & 70B & 35.2 & 56.8 & 13.5\\
% \gmidrule
% \bf Mistral & 7B & 29.4 & 47.5 & 11.3\\
% \gmidrule
% \bf InternLM & 20B & 36.9 & 62.9 & 10.9 \\
% \gmidrule
% \bf Skywork & 7B & 31.8 & 55.8 & 7.8 \\
% \gmidrule
% \multirow{2}{*}{\bf Yi}  
% & 6B & 18.6 & 32.5 & 4.6 \\
% & 34B & 40.8 & 67.2 & 14.4 \\
% \bottomrule

%   \end{tabular}
%    \caption{Comparison between models on GSM8k and MATH.
%    % \alex{what kind of prompt did we use? do we test the numbers of all these models on our own, or we copy the number from their papers?}
%    }
%   \label{tab:math}
% \end{table}

% \begin{table}
%   \centering
%   \begin{tabular}{lcccccc}
% \toprule
%               \bf Model & \bf Size & \bf Avg & \bf GSM8k & \bf MATH & \bf HumanEval & \bf MBPP\\
% \midrule
% \bf GPT-3.5 & - & 35.6 & 57.1 & 14.0 & 48.1 & 61.4\\
% \bf GPT-4 & - & \textbf{66.1} & \textbf{92} & \textbf{40.2} & 67.0 & 63.6 \\
% \gmidrule
% \bf Falcon & 180B & - & 54.4 & - & 0.61 & 47.0\\
% \gmidrule
% \multirow{2}{*}{\bf Qwen}
%  & 7B & 31.7 & 51.7 & 11.6 & 29.9 & 34.0\\
% & 14B & 43.1 & 61.3 & 24.8 & 32.3 & 48.9 \\
%  \gmidrule
% \multirow{2}{*}{\bf Baichuan 2}
% & 7B & 15.1 & 24.5 & 5.6 \\
% & 13B & 16.1 & 22.1 & 10.1 \\
% % \midrule
% % Open-L\textsmaller{LaMA} v2 & 7B & 9.7 & 2.1 \\
% \gmidrule
% \multirow{3}{*}{\bf L\textsmaller{LaMA} 2} 
% & 7B & 10.0 & 16.7 & 3.3\\ 
% & 34B & 24.2 & 42.2 & 6.2\\ 
% & 70B & 35.2 & 56.8 & 13.5\\
% \gmidrule
% \bf Mistral & 7B & 29.4 & 47.5 & 11.3\\
% \gmidrule
% \bf InternLM & 20B & 36.9 & 62.9 & 10.9 \\
% \gmidrule
% \bf Skywork & 7B & 31.8 & 55.8 & 7.8 \\
% \gmidrule
% \multirow{2}{*}{\bf Yi}  
% & 6B & 18.6 & 32.5 & 4.6 \\
% & 34B & 40.8 & 67.2 & 14.4 \\
% \bottomrule

%   \end{tabular}
%    \caption{Comparison between models on GSM8k and MATH.
%    % \alex{what kind of prompt did we use? do we test the numbers of all these models on our own, or we copy the number from their papers?}
%    }
%   \label{tab:math}
% \end{table}


\begin{table}
  \centering
  \begin{tabular}{@{}lccccc@{}}
    \toprule
    \bf Model & \bf Size & \bf GSM8k & \bf MATH & \bf Human-Eval pass@1 & \bf MBPP pass@1 \\
    \midrule
    \bf GPT-3.5 & - & 57.1 & 14.0 & 48.1 & 61.4 \\
    \bf GPT-4 & - & \textbf{92.0} & \textbf{40.2} & \textbf{67.0} & \textbf{63.6} \\
    \gmidrule
    \bf Falcon & 180B  & 54.4 & - & 0.61 & 47.0 \\
    \gmidrule
    \multirow{2}{*}{\bf Qwen}
    & 7B  & 51.7 & 11.6 & 29.9 & 34.0 \\
    & 14B  & 61.3 & 24.8 & 32.3 & 48.9 \\
    \gmidrule
    \multirow{2}{*}{\bf Baichuan 2}
    & 7B & 24.5 & 5.6 & 18.3 & 28.3 \\
    & 13B & 22.1 & 10.1 & 20.7 & 26.1 \\
    \gmidrule
    \multirow{3}{*}{\bf L\textsmaller{LaMA} 2} 
    & 7B  & 16.7 & 3.3 & 12.8 & 14.8\\ 
    & 34B  & 42.2 & 6.2 & 22.6 & 33.0\\ 
    & 70B  & 56.8 & 13.5 & 31.7 & 45.0\\
    \gmidrule
    \bf Mistral & 7B & 47.5 & 11.3 & 30.5 & 47.5 \\
    \gmidrule
    \bf InternLM & 20B & 62.9 & 10.9 & 28.1 & 41.4 \\
    \gmidrule
    \bf Skywork & 7B & 55.8 & 7.8 & 13.4 & 22.8 \\
    \gmidrule
    \multirow{2}{*}{\bf Yi}  
    & 6B & 32.5 & 4.6 & 15.9 & 26.3 \\
    & 34B & 67.2 & 14.4 & 23.2 & 41.0 \\
    \bottomrule
  \end{tabular}
  \caption{Comparison of models on GSM8k, MATH, Human-Eval, and MBPP.}
  \label{tab:code_math}
\end{table}



By training on a significantly larger number of tokens (3.1T) compared to prior work (usually $\le$ 2T), we have observed a substantial performance gain across  benchmarks, as shown in Table \ref{tab:models}. However, it is important to note that there are still discernible disparities between our model and existing open-source and close-source models, particularly in tasks related to mathematics and coding. As performance in these domains can be significantly improved by continual pretraining and instruction fine-tuning, we have refrained from incorporating extensive mathematical and coding content in the pretraining corpus when making the initial design choices. We do plan to release models with enhanced math and coding capabilities in the future.


\begin{figure}
	\begin{center}
		\includegraphics[width=\textwidth]{images/icl.pdf}
	\end{center}
	\caption{Evaluating language model's in-context learning capability by inferring the linear coefficients of a weighted sum.
        Considering the discussions of whether emergent ability is an artifact of measurement~\citep{schaeffer2024emergent}, we use 
        difference to the target (target number - model prediction) as a continuous measure, 
        and exact match (target number == model prediction) as a discontinuous measure.
        A: when there is two linear coefficients, Yi-34B performs the best when measuring by the difference to the target number. 
        B: increasing the number of linear coefficients to 5, only models that are large enough (LLaMA2 70B and Mixtral 8x7B) can achieve meaningful exact match, showing that in-context learning complex functions is an emergent ability.
        } %
	\label{fig:icl}
\end{figure}



\subsubsection{Discussions}
% We discuss several eye-catching observations in the section and give empirical conclusions for assisting future LLM pretrain research works.
% There are a few important takes from our experiment results:

\begin{description}[leftmargin=0.2in]
\item[Gain from Model Scale.] We observe that Yi-34B has substantial performance improvement compared to Yi-6B, though they utilized the same pretrain corpora.
Larger model size leads to higher performance gain on Code and Math benchmarks, referring to Tab.~\ref{tab:code_math}, compared to benchmarks focusing on Commonsense Reasoning, Reading Comprehension, or Knowledge.
\item[Data Quality.]
Smaller models of higher quality pretrain data, like Yi-34B or Qwen-14B, usually demonstrate better performance than models of larger size but (presumably) lower quality data, such as Falcon-180B (though the focus of Falcon-180B might be more on the scaling side, which is definitely of important value on its own). 
\item[Gap between GPT-4 and Open-source LLMs.]
Based on Tab.~\ref{tab:models}, we note that open-source LLMs still lag behind the performance of GPT-4 and GPT-3.5 on various benchmarks.
Yet representative bilingual LLMs, e.g. Qwen-14B and Yi-34B, can match or even surpass the performance of GPT-4 on Chinese knowledge related benchmarks, including C-Eval~\cite{huang2023c}, CMMLU~\cite{li2023cmmlu}, and Gaokao~\cite{zhang2023evaluating}.
However, there is still a huge gap between GPT-4 and open-source models on reasoning-related benchmarks like BBH~\cite{srivastava2023imitation}, code (HumanEval), and math (MATH).
\end{description}

\subsubsection{In-Context Learning Study}
% \begin{figure}
	\begin{center}
		\includegraphics[width=\textwidth]{images/icl.pdf}
	\end{center}
	\caption{Evaluating language model's in-context learning capability by inferring the linear coefficients of a weighted sum.
        Considering the discussions of whether emergent ability is an artifact of measurement~\citep{schaeffer2024emergent}, we use 
        difference to the target (target number - model prediction) as a continuous measure, 
        and exact match (target number == model prediction) as a discontinuous measure.
        A: when there is two linear coefficients, Yi-34B performs the best when measuring by the difference to the target number. 
        B: increasing the number of linear coefficients to 5, only models that are large enough (LLaMA2 70B and Mixtral 8x7B) can achieve meaningful exact match, showing that in-context learning complex functions is an emergent ability.
        } %
	\label{fig:icl}
\end{figure}
We further investigate the in-context learning capability, i.e., the capability of inferring the underlying function given the few-show input-output demonstrations.
We consider the task of inferring the linear coefficient of a weighted sum.
Specifically, define $y = w_1x_1 + w2x_2 + ... + w_nx_n$, our few-shot demonstration is $x_1, x_2, ..., x_n, y$, and we ask the model to (implicitly) infer $w_1, w_2, ..., w_n$ by predicting the $y$ given a new set of input $x$. 
We use (a). the absolute difference between model prediction $y$ and the ground truth $y^*$, i.e., $|y - y^*|$ as a continuous measure, and use (b). the exact match $y == y^*$ as a discontinuous measure.
We further note that most of the models perform reasonably well on addition and subtraction, so the ability to do arithmetic, as a confounding factor, can be ruled out. 

The results are shown in Figure~\ref{fig:icl}.
When setting the linear coefficients of be [1, -1], we see that Yi 34B and LLaMA-2 70B performs the best in-terms of answer exact match. 
If we increase the number of the linear coefficients to be [1, 1, 1, 1, 1], we observe the emergent behavior that only large models (LLaMA-2 70B and Mixtral) can achieve good scores on exact match, although the differences to target is more continuous. 
These observations give side evidence for Yi-34B's performance on in-context learning and indicates that further scaling may allow the model to infer more complicated functions by in-context learning.